\section{Period and lattice certificates}
\label{app:A}

This appendix records the explicit witnesses for parts~2 and~3 of
Proposition~\ref{prop:graph}. All strings below are the sequences of
\emph{emitted letters} of a closed walk read from the stated root; the walk
returns to that root. Every walk was checked edge by edge against the
memory-\(m\) window graph, and the state validity of the root was checked
independently.

For each \(h\) the period certificate consists of two closed walks at the root
whose lengths are coprime, which forces graph period \(1\). The lattice
certificate consists of three closed walks of one common length \(L\) at the
same root; writing \(\Psi\) for the Parikh vector and projecting onto the first
two coordinates, the two differences
\(d_1=\Psi(c_2)-\Psi(c_1)\) and \(d_2=\Psi(c_3)-\Psi(c_1)\) satisfy
\(|\det[d_1\ d_2]|=1\), so they generate \(\mathbb Z^2\) and the projected
lattice has covolume \(1\).

\subsection*{\texorpdfstring{\(h=2\)}{h=2}, \texorpdfstring{\(m=3\)}{m=3}, root \texttt{aaa}}
\begin{itemize}
\item Period: closed walks of lengths \(4\) (\texttt{baaa}) and \(5\)
      (\texttt{bcaaa}); \(\gcd(4,5)=1\).
\item Lattice at \(L=6\):
      \(c_1=\texttt{bacaaa}\), \(\Psi=(4,1,1)\);
      \(c_2=\texttt{bbbaaa}\), \(\Psi=(3,3,0)\);
      \(c_3=\texttt{bbcaaa}\), \(\Psi=(3,2,1)\).
\item \(d_1=(-1,2)\), \(d_2=(-1,1)\), \(\det=1\).
\end{itemize}

\subsection*{\texorpdfstring{\(h=3\)}{h=3}, \texorpdfstring{\(m=5\)}{m=5}, root \texttt{aabaa}}
\begin{itemize}
\item Period: closed walks of lengths \(6\) (\texttt{caabaa}) and \(7\)
      (\texttt{acaabaa}); \(\gcd(6,7)=1\).
\item Lattice at \(L=9\):
      \(c_1=\texttt{acbcaabaa}\), \(\Psi=(5,2,2)\);
      \(c_2=\texttt{acccaabaa}\), \(\Psi=(5,1,3)\);
      \(c_3=\texttt{cbccaabaa}\), \(\Psi=(4,2,3)\).
\item \(d_1=(0,-1)\), \(d_2=(-1,0)\), \(\det=-1\).
\end{itemize}

\subsection*{\texorpdfstring{\(h=4\)}{h=4}, \texorpdfstring{\(m=7\)}{m=7}, root \texttt{aaabaaa}}
\begin{itemize}
\item Period: closed walks of lengths \(8\) (\texttt{caaabaaa}) and \(11\)
      (\texttt{cabcaaabaaa}); \(\gcd(8,11)=1\).
\item Lattice at \(L=13\):
      \(c_1=\texttt{caaabcaaabaaa}\), \(\Psi=(9,2,2)\);
      \(c_2=\texttt{caabccaaabaaa}\), \(\Psi=(8,2,3)\);
      \(c_3=\texttt{cabbccaaabaaa}\), \(\Psi=(7,3,3)\).
\item \(d_1=(-1,0)\), \(d_2=(-2,1)\), \(\det=-1\).
\end{itemize}

\subsection*{\texorpdfstring{\(h=5\)}{h=5}, \texorpdfstring{\(m=9\)}{m=9}, root \texttt{aaabaaaca}}
\begin{itemize}
\item Period: closed walks of lengths \(8\) (\texttt{aabaaaca}) and \(13\)
      (\texttt{abccaaabaaaca}); \(\gcd(8,13)=1\).
\item Lattice at \(L=13\):
      \(c_1=\texttt{abccaaabaaaca}\), \(\Psi=(8,2,3)\);
      \(c_2=\texttt{bbccaaabaaaca}\), \(\Psi=(7,3,3)\);
      \(c_3=\texttt{bcccaaabaaaca}\), \(\Psi=(7,2,4)\).
\item \(d_1=(-1,1)\), \(d_2=(-1,0)\), \(\det=1\).
\end{itemize}

\subsection*{\texorpdfstring{\(h=6\)}{h=6}, \texorpdfstring{\(m=11\)}{m=11}, root \texttt{abbbcaaaccb}}
\begin{itemize}
\item Period: closed walks of lengths \(12\) (\texttt{babbbcaaaccb}) and \(17\)
      (\texttt{acccaaabbbcaaaccb}); \(\gcd(12,17)=1\).
\item Lattice at \(L=17\):
      \(c_1=\texttt{acccbbabbbcaaaccb}\), \(\Psi=(5,6,6)\);
      \(c_2=\texttt{baaacaabbbcaaaccb}\), \(\Psi=(8,5,4)\);
      \(c_3=\texttt{bbccaaabbbcaaaccb}\), \(\Psi=(6,6,5)\).
\item \(d_1=(3,-1)\), \(d_2=(1,0)\), \(\det=1\).
\end{itemize}

\subsection*{Root independence}

The projected lattice generated by differences of equal-length closed walks
does not depend on the root chosen. If \(c,c'\) are equal-length closed walks
at \(z\), and \(p\) is a path \(z'\to z\) and \(q\) a path \(z\to z'\) inside
the same strongly connected component, then \(pcq\) and \(pc'q\) are
equal-length closed walks at \(z'\) with the same Parikh difference. A single
root therefore suffices for each \(h\).
